% --- Archon physics formalization source begin ---
% archon:physics
% archon:covers IPhO2026Problems/problem_IPhO_2026_2_B_2.lean
% archon:source-report reports/ipho_2026_k3/problem_IPhO_2026_2_B_2.source.json
% archon:problem-id IPhO_2026_2
% archon:part-id B.2
% archon:previous-part IPhO_2026_2_B_1
% archon:previous-part-policy natural_language_prerequisite_only

\chapter{Physics problem IPhO\_2026\_2\_B\_2}
\label{ch:IPhO2026Problems_problem_IPhO_2026_2_B_2}

\paragraph{Problem source.}
A half-cylindrical mirror of radius R illuminates a fully absorbing cylindrical
container of radius a.  Their axes are parallel, and the container center lies
R/2 from the mirror center on the symmetry plane.  Uniform parallel sunlight
arrives along the optical axis.  Any ray absorbed by the container reflects at
most once.  Let theta\_max be the largest incidence angle on the mirror among
rays that strike the container, and let P\_0 be the power the cylinder would
receive without the mirror.  See Figure 2f.

Current subquestion:
Express the power ratio P/P\_0 in terms of theta\_max.

\paragraph{Current subquestion.}
Express the power ratio P/P\_0 in terms of theta\_max.

\paragraph{Recorded answer/context.}
P/P\_0 = 1/(1 - cos(theta\_max)).

\paragraph{Figure/image path.}
/root/proposal\_for\_physic/science-mango/ipho\_2026\_source/image/T2\_page-3.png

\paragraph{Reusable previous-part conclusions.}
\begin{itemize}
\item Source B.1. Question: Given a = alpha*sin(theta\_max) + beta*sin(2*theta\_max), determine alpha and beta in terms of R. Reusable conclusions: alpha = R and beta = -R/2. Policy: natural\_language\_prerequisite\_only; do\_not\_import\_Lean\_output
\end{itemize}

\paragraph{Formalization target.}
create a compiling Lean file with sorry bodies at `IPhO2026Problems/problem\_IPhO\_2026\_2\_B\_2.lean`.
The Lean declarations must preserve the physical quantities, dimensions or dimensional roles, figure labels, governing-law hypotheses, and final relation expressed by this problem.
Use Mathlib/Physlib names found through LeanExplore where available. If a domain API is missing, introduce faithful local abstractions rather than scalar placeholder aliases.

\begin{theorem}[Physics formalization target]
\label{thm:physics:IPhO_2026_2_B_2:target}
\uses{thm:IPhO2026Problems_problem_IPhO_2026_2_B_2:power_ratio_in_terms_of_theta_max}
The assigned autoformalize agent should translate this physics problem into Lean declarations in the covered file, with theorem and lemma proof bodies written as `by sorry`.
\end{theorem}
\begin{proof}
This is an autoformalization task, not a proof task. Produce faithful statements that can later be proved without weakening the source contract.
\end{proof}
% --- Archon physics formalization source end ---

% --- Archon named-quantities coverage (blueprint-writer 2-b-2-entries; restored planner-side iter-009 after the chapter body was found truncated back to the 51-line iter-002 skeleton) ---
% NOTE: import policy (mirrors the iter-002 PhysLean-coverage exemption NOTE of
% this chapter family, e.g. 2-b-1): the covered file builds on the
% `import Mathlib` baseline alone; PhysLean has no specular-reflection /
% geometric-optics module for the Figure-2f half-cylinder mirror regime, so no
% Physlib import is added (see the physics-grounding log for this file).  This
% NOTE records the resolved import policy for the blueprint-doctor
% `missing-physlib-import` check.

\subsection*{Cross-section geometry and data}

\begin{definition}[Cross-section plane]
\label{def:IPhO2026Problems_problem_IPhO_2026_2_B_2:Plane}
\lean{IPhO2026_2_B_2.Plane}
The transverse cross-sectional plane of the cooker, the Euclidean plane
$\mathbb{E}^2 \cong \mathbb{R}^2$ with its inner product.  The physical
system is translationally invariant along the cylinder axes, so a single
cross-section captures all the geometry; points have units of length.
\end{definition}
\begin{proof}
Definition; no claim.
\end{proof}

\begin{definition}[Cooker parameters]
\label{def:IPhO2026Problems_problem_IPhO_2026_2_B_2:CookerParams}
\lean{IPhO2026_2_B_2.CookerParams}
Dimensionful parameters of the cooker: the mirror radius $R$ and the
container radius $a$ (both lengths), bundled with their positivity
certificates $0 < R$ and $0 < a$.
\end{definition}
\begin{proof}
Definition; no claim.
\end{proof}

\begin{definition}[Cooker cross-section geometry]
\label{def:IPhO2026Problems_problem_IPhO_2026_2_B_2:CookerGeometry}
\lean{IPhO2026_2_B_2.CookerGeometry}
\uses{def:IPhO2026Problems_problem_IPhO_2026_2_B_2:CookerParams, def:IPhO2026Problems_problem_IPhO_2026_2_B_2:Plane}
The geometry of the cross-section (Figure~2f): the mirror centre $C$, the
container centre $A$, the unit vector $e$ along the optical axis pointing
from $C$ towards the open half of the half-cylinder mirror, and the
in-plane unit normal $n$ ($n \perp e$, chosen on the container's side),
with the container centre lying $R/2$ from $C$ along $n$:
$A - C = (R/2) \bullet n$.  Sunlight arrives in the direction $-e$.
The bundled fields record $\|e\| = 1$, $\|n\| = 1$,
$\langle n, e \rangle = 0$, and the offset law.
\end{definition}
\begin{proof}
Definition; no claim.
\end{proof}

\begin{definition}[Mirror circle]
\label{def:IPhO2026Problems_problem_IPhO_2026_2_B_2:mirrorCircle}
\lean{IPhO2026_2_B_2.mirrorCircle}
\uses{def:IPhO2026Problems_problem_IPhO_2026_2_B_2:CookerParams, def:IPhO2026Problems_problem_IPhO_2026_2_B_2:CookerGeometry}
The full circle of radius $R$ centred at $C$: the cross-sectional profile
of the cylindrical mirror.
\end{definition}
\begin{proof}
Definition; no claim.
\end{proof}

\begin{definition}[Container disk]
\label{def:IPhO2026Problems_problem_IPhO_2026_2_B_2:containerDisk}
\lean{IPhO2026_2_B_2.containerDisk}
\uses{def:IPhO2026Problems_problem_IPhO_2026_2_B_2:CookerParams, def:IPhO2026Problems_problem_IPhO_2026_2_B_2:CookerGeometry}
The absorbing disc: the closed disk of radius $a$ centred at $A$, the
cross-section of the cylindrical container.
\end{definition}
\begin{proof}
Definition; no claim.
\end{proof}

\begin{definition}[Half-mirror arc]
\label{def:IPhO2026Problems_problem_IPhO_2026_2_B_2:halfMirrorArc}
\lean{IPhO2026_2_B_2.halfMirrorArc}
\uses{def:IPhO2026Problems_problem_IPhO_2026_2_B_2:mirrorCircle}
The half-cylinder mirror arc: the semicircle of
\cref{def:IPhO2026Problems_problem_IPhO_2026_2_B_2:mirrorCircle} lying in
the $0 \le \langle m - C, e\rangle$ half-plane --- the
``half-hollow-cylinder mirror'' of Figure~2f.
\end{definition}
\begin{proof}
Definition; no claim.
\end{proof}

\subsection*{Governing-law structures}

\begin{definition}[Absorbed-ray bookkeeping]
\label{def:IPhO2026Problems_problem_IPhO_2026_2_B_2:AbsorbedRays}
\lean{IPhO2026_2_B_2.AbsorbedRays}
\uses{def:IPhO2026Problems_problem_IPhO_2026_2_B_2:halfMirrorArc, def:IPhO2026Problems_problem_IPhO_2026_2_B_2:containerDisk, def:IPhO2026Problems_problem_IPhO_2026_2_B_2:CookerParams}
Single-bounce bookkeeping for the sun rays absorbed by the container.  The
carrier pairs the incidence-point map $\mathrm{incidentPt} : \mathbb{R} \to
\mathbb{E}^2$ with the impact-parameter hit set; the bundled hypotheses
record the governing physical laws, kept as assumptions rather than
redefined locally: every absorbed ray (a) strikes the half-cylinder arc of
the mirror (incidence points land in
\cref{def:IPhO2026Problems_problem_IPhO_2026_2_B_2:halfMirrorArc}), (b)
obeys the specular reflection law on the circular profile --- the reflected
ray reaches an absorbed endpoint $q$ in
\cref{def:IPhO2026Problems_problem_IPhO_2026_2_B_2:containerDisk} with
$q - \mathrm{incidentPt}\,y = (2\langle \mathrm{incidentPt}\,y - C, n\rangle - R) \bullet n$ ---
and the Figure-2f branch information (c) the axial (central) ray is
absorbed, (d) the absorbed family is gap-free right of the axis
(contiguity of the single-bounce fan), and (e) every impact parameter in
$(0, R)$ is realized (strengthened at iter-009, see the NOTE below), so
the collected fan is exactly the open half aperture: it fills one side
from the axis to the aperture edge and $\theta_{\max}$ is a maximum over
a connected family.
\end{definition}
\begin{proof}
Definition; no claim.
\end{proof}

% NOTE: Statement reconciliation (planner-recorded, iter-009; review-gate
% session-8 R1).  Clause (e) previously read ``every impact parameter in
% $(0, a)$ is realized'', quantifying only over the container silhouette
% band.  That form cannot entail the advertised
% $\mathrm{collectedWidth} = R$: at container radius $a \ll R$ a thin
% absorbed fan inside $(0, a)$ satisfies every field while the supremum
% stays $\le a$ (constructive countermodel at $R = 1$, $a = 0.1$:
% $P/P_0 = R/(2a) = 5 \neq 1/(1-\cos\theta_{\max}) \approx 1.005$).
% The field was therefore strengthened to quantify over $(0, R)$, the full
% Figure-2f half aperture the proof of
% \cref{lem:IPhO2026Problems_problem_IPhO_2026_2_B_2:collectedWidth_eq_radius}
% already narrates.  Source warrant: under the B.1 calibration
% $a = R\sin\theta\,(1-\cos\theta)$, $\theta \in (0, \/2)$, a
% figure-realizable family at arbitrarily small $a$ exists only if coverage
% reaches arbitrarily close to $R$ --- its $\theta$-value must also realize
% the $\theta_{\max}$ specification of
% \cref{def:IPhO2026Problems_problem_IPhO_2026_2_B_2:ThetaMaxSpec}.  With
% the strengthened field the $\mathrm{collectedWidth} = R$ proof closes
% (supremum over the open half aperture); without it the target theorem is
% false.
% End of NOTE.

\begin{definition}[Collected width]
\label{def:IPhO2026Problems_problem_IPhO_2026_2_B_2:collectedWidth}
\lean{IPhO2026_2_B_2.collectedWidth}
\uses{def:IPhO2026Problems_problem_IPhO_2026_2_B_2:AbsorbedRays}
The transverse width (impact-parameter support) of the rays that the mirror
concentrates onto the container: the supremum of the axial readouts
$\langle \mathrm{incidentPt}\,y - C, e\rangle$ over the hit set.  With
uniform parallel sunlight this width is proportional to the received power
$P$ (per unit axial length).  At the Figure-2f operating point the
hit-set readouts fill the open half aperture $(0, R)$ (full side
coverage), so the supremum is $R$ --- the physical content of
\cref{lem:IPhO2026Problems_problem_IPhO_2026_2_B_2:collectedWidth_eq_radius}.
\end{definition}
\begin{proof}
Definition; no claim.
\end{proof}

\begin{definition}[Uniform sunlight intensity]
\label{def:IPhO2026Problems_problem_IPhO_2026_2_B_2:UniformIntensity}
\lean{IPhO2026_2_B_2.UniformIntensity}
Uniform parallel sunlight: the incoming intensity $I$ (power per unit area)
is constant across the aperture and certified positive by a bundled field
$0 < I$.
\end{definition}
\begin{proof}
Definition; no claim.
\end{proof}

\begin{definition}[Power budget]
\label{def:IPhO2026Problems_problem_IPhO_2026_2_B_2:PowerBudget}
\lean{IPhO2026_2_B_2.PowerBudget}
\uses{def:IPhO2026Problems_problem_IPhO_2026_2_B_2:collectedWidth, def:IPhO2026Problems_problem_IPhO_2026_2_B_2:UniformIntensity, def:IPhO2026Problems_problem_IPhO_2026_2_B_2:CookerParams}
The power budget for the configuration: $P$, the power actually received by
the absorbing container from the mirror, and $P_0$, the power it would
receive without the mirror.  The bundled hypotheses state the physical
accounting under uniform intensity $I$: both quantities are proportional to
their transverse collected widths per unit axial length, namely
$P = I \cdot \mathrm{collectedWidth}$ and $P_0 = I \cdot (2a)$ (the
container's unmirrored geometric width, its diameter).
\end{definition}
\begin{proof}
Definition; no claim.
\end{proof}

\begin{definition}[B.1 calibration]
\label{def:IPhO2026Problems_problem_IPhO_2026_2_B_2:B1Calibration}
\lean{IPhO2026_2_B_2.B1Calibration}
\uses{def:IPhO2026Problems_problem_IPhO_2026_2_B_2:CookerParams}
The previous-part (B.1, $\alpha = R$, $\beta = -R/2$) result, kept as a
calibration hypothesis for the geometry: the container radius relates to
the maximum incidence angle by
$a = R \sin\theta - (R/2) \sin(2\theta)$.
\end{definition}
\begin{proof}
Definition; no claim.
\end{proof}

\begin{definition}[Incidence angle]
\label{def:IPhO2026Problems_problem_IPhO_2026_2_B_2:incidenceAngle}
\lean{IPhO2026_2_B_2.incidenceAngle}
\uses{def:IPhO2026Problems_problem_IPhO_2026_2_B_2:CookerParams, def:IPhO2026Problems_problem_IPhO_2026_2_B_2:CookerGeometry}
The incidence angle at a mirror point $m$ on the circular profile of
radius $R$, measured against the normal at $m$:
$\arccos\bigl(|\langle m - C, e\rangle| / R\bigr)$ --- the angle
between the incoming direction $e$ and the radius to $m$.
\end{definition}
\begin{proof}
Definition; no claim.
\end{proof}

\begin{definition}[Maximum incidence angle specification]
\label{def:IPhO2026Problems_problem_IPhO_2026_2_B_2:ThetaMaxSpec}
\lean{IPhO2026_2_B_2.ThetaMaxSpec}
\uses{def:IPhO2026Problems_problem_IPhO_2026_2_B_2:AbsorbedRays, def:IPhO2026Problems_problem_IPhO_2026_2_B_2:incidenceAngle}
The specification of $\theta_{\max}$: the largest angle of incidence on
the mirror (measured against the normal at the point of incidence) among
all reflected rays striking the container, lying in $(0, \/2)$ --- the
branch of nontrivial, non-grazing rays seen in Figure~2f.  Concretely:
$\theta$ is attained at some absorbed impact parameter and bounds the
incidence angle of every absorbed ray.
\end{definition}
\begin{proof}
Definition; no claim.
\end{proof}

\subsection*{Derivation bridges}

\begin{lemma}[Impact parameters bounded by the aperture]
\label{lem:IPhO2026Problems_problem_IPhO_2026_2_B_2:impactParam_le_aperture}
\lean{IPhO2026_2_B_2.impactParam_le_aperture}
\uses{def:IPhO2026Problems_problem_IPhO_2026_2_B_2:AbsorbedRays, def:IPhO2026Problems_problem_IPhO_2026_2_B_2:mirrorCircle}
Every mirror-collected impact parameter is bounded in absolute value by
the aperture radius: for every absorbed impact parameter $y$,
$|\langle \mathrm{incidentPt}\,y - C, e\rangle| \le R$.
\end{lemma}
\begin{proof}
The incidence point lies on the mirror circle, so
$\|\mathrm{incidentPt}\,y - C\| = R$; since $\|e\| = 1$,
Cauchy--Schwarz for the inner product gives
$|\langle \mathrm{incidentPt}\,y - C, e\rangle| \le
\|\mathrm{incidentPt}\,y - C\| \cdot \|e\| = R$.
\end{proof}

\begin{lemma}[Collected width equals the radius]
\label{lem:IPhO2026Problems_problem_IPhO_2026_2_B_2:collectedWidth_eq_radius}
\lean{IPhO2026_2_B_2.collectedWidth_eq_radius}
\uses{def:IPhO2026Problems_problem_IPhO_2026_2_B_2:collectedWidth, def:IPhO2026Problems_problem_IPhO_2026_2_B_2:AbsorbedRays, lem:IPhO2026Problems_problem_IPhO_2026_2_B_2:impactParam_le_aperture, def:IPhO2026Problems_problem_IPhO_2026_2_B_2:CookerParams}
The transverse width collected by the half-cylinder mirror on the
container's side equals $R$ (one full half of the aperture $2R$):
$\mathrm{collectedWidth} = R$.
\end{lemma}
\begin{proof}
Upper bound: every impact-parameter readout is at most $R$ by
\cref{lem:IPhO2026Problems_problem_IPhO_2026_2_B_2:impactParam_le_aperture}
(incidence points lie on the mirror circle; Cauchy--Schwarz with
$\|e\| = 1$), so the defining supremum does not exceed $R$; the readout
family is bounded and nonempty (the central ray $0$ is absorbed).
Lower bound: the (iter-009 strengthened) full side coverage realizes every
impact parameter $y \in (0, R)$ as an absorbed readout, so the readout
set contains the whole open half aperture $(0, R)$, whose supremum in
$\mathbb{R}$ is $R$.  Hence the supremum of the full readout family is at
least $R$; combined with the upper bound it equals $R$.
\end{proof}

\begin{lemma}[Power ratio as width ratio]
\label{lem:IPhO2026Problems_problem_IPhO_2026_2_B_2:power_ratio_eq_width_ratio}
\lean{IPhO2026_2_B_2.power_ratio_eq_width_ratio}
\uses{def:IPhO2026Problems_problem_IPhO_2026_2_B_2:PowerBudget, lem:IPhO2026Problems_problem_IPhO_2026_2_B_2:collectedWidth_eq_radius}
Power ratio from the width accounting only:
$P / P_0 = R / (2a)$.  With the B.1 calibration this becomes the target
trigonometric form.
\end{lemma}
\begin{proof}
Substitute the two budget identities
$P = I \cdot \mathrm{collectedWidth}$ and $P_0 = I \cdot (2a)$; the
common positive intensity $I$ cancels in the ratio, and
\cref{lem:IPhO2026Problems_problem_IPhO_2026_2_B_2:collectedWidth_eq_radius}
replaces $\mathrm{collectedWidth}$ by $R$.
\end{proof}

\begin{lemma}[Radius-over-diameter trigonometric bridge]
\label{lem:IPhO2026Problems_problem_IPhO_2026_2_B_2:radius_over_diameter_eq}
\lean{IPhO2026_2_B_2.radius_over_diameter_eq}
\uses{def:IPhO2026Problems_problem_IPhO_2026_2_B_2:B1Calibration}
For $\theta \in (0, \/2)$ satisfying the B.1 calibration,
$R / (2a) = 1 / (1 - \cos\theta)$.
\end{lemma}
\begin{proof}
By the double-angle identity and the calibration,
$2a = 2R\sin\theta - R\sin(2\theta)
= 2R\sin\theta\,(1 - \cos\theta)$,
so $R/(2a) = 1/(1 - \cos\theta)$ after cancelling the positive factor
$2R\sin\theta$ (positivity from $\theta \in (0, \/2)$ and $0 < R$);
this leaves $R/(2a) = 1/(1 - \cos\theta)$.
\end{proof}

\begin{theorem}[Power ratio in terms of $\theta_{\max}$ (T2-B.2 target)]
\label{thm:IPhO2026Problems_problem_IPhO_2026_2_B_2:power_ratio_in_terms_of_theta_max}
\lean{IPhO2026_2_B_2.power_ratio_in_terms_of_theta_max}
\uses{lem:IPhO2026Problems_problem_IPhO_2026_2_B_2:power_ratio_eq_width_ratio, lem:IPhO2026Problems_problem_IPhO_2026_2_B_2:radius_over_diameter_eq, def:IPhO2026Problems_problem_IPhO_2026_2_B_2:ThetaMaxSpec}
For the cooker of Figure~2f with absorbed-ray bookkeeping, power budget,
and maximum-incidence-angle specification $\theta_{\max}$, the ratio of
the received power $P$ to the unmirrored power $P_0$ is
$P / P_0 = 1 / (1 - \cos\theta_{\max})$ --- the recorded official
answer of part B.2.
\end{theorem}
\begin{proof}
The B.1 calibration at the specified maximum angle yields
$R/(2a) = 1/(1 - \cos\theta_{\max})$ by
\cref{lem:IPhO2026Problems_problem_IPhO_2026_2_B_2:radius_over_diameter_eq};
chaining with
\cref{lem:IPhO2026Problems_problem_IPhO_2026_2_B_2:power_ratio_eq_width_ratio},
$P/P_0 = R/(2a)$, yields the claim.
\end{proof}
